% Transcription — fonds Alexandre Grothendieck, shelfmark 155, batch 1 (pages 1-10).
% Established from the facsimile archives/batches/155/batch-01.pdf (a mirror of
% https://grothendieck.umontpellier.fr/155.pdf), per the skill
% transcribe-grothendieck. The folder is ten pages and this is its only batch.
%
% Five of the ten leaves carry mathematics in his hand — the odd-numbered ones,
% 1, 3, 5, 7, 9 — and they carry his own pagination, a circled 1 to 5 in the
% top right corner, so the run is continuous and complete. The five even leaves
% are other people's paper (a publishing contract, three pieces of USTL and
% Paul Valéry administration, a colleague's circular) with nothing of his on
% them; they are skipped, and the gaps in the page numbers are the record.
%
% The run is one argument and it is the one the archivists' bracketed title
% guesses at. Page 1 states the dictionary: over a finite extension k'/k, a
% tensor category with an exact faithful ACU fibre functor to Mod_f(k') is the
% same thing as an affine scheme G' over k' with a section — a commutative
% k'-algebra U with augmentation epsilon — together with, on the pro-object
% A = U-dual, a commutative associative structure with epsilon as unit, subject
% to two multiplicativity conditions, the second of which takes its values in
% the subquotient of A (completed tensor over k) A cut out by the ideal
% generated by J = Ker(k' tensor_k k' --> k'). Page 3 runs the example
% C = Mod_f(k) with F = - tensor_k k', where U = k' tensor_k k' and
% A = End_k(k'), and then asks the question of the folder: what conditions on C
% — intrinsic ones — make the base-change functor exact and faithful, and does
% it suffice that C be rigid with k = End(1_C)? « Exprimer la rigidité de C ! »
% Pages 5 to 9 are the attempt: rigidity as an exact equivalence M |--> M-dual
% compatible with F', represented by Hom_A(M, I) with I = D(A) = U, the two
% A-module structures on U, End_A(U) = A-opposite and the involution
% a |--> a-check it forces, then the computation of F'(D(M)) = M-dual over k'
% and the compatibility the evaluation pairing must satisfy. The last page
% reduces that compatibility to the case M = A and breaks off on an unfinished
% « A tensor ».
%
% The hand is drafting at speed. The formulas, the underlined words and the
% displayed maps carry; a good deal of the connective prose does not, and is
% marked rather than filled. Two blocks run diagonally up the left edge of
% pages 3 and 5 in a hand faster still and are given only as far as they go.
% Three readings are flagged and not settled: « ACU » for the abbreviation that
% qualifies both the category and the functor on page 1 (associative,
% commutative, unit — the word is nowhere written out here); the structure word
% of article 2°) on the same page, under a strike; and the two J-glyphs of the
% same article, which the leaf distinguishes only by a tilde.
%
% Pass: Opus 5 (claude-opus-5), 2026-09-12 — first pass, unchecked against the pages by a human.
\documentclass[11pt,a4paper]{article}
% Le préambule commun à toutes les transcriptions du fonds Grothendieck.
%
% Il tient en une idée : **l'incertitude se compose, elle ne se résout pas.**
% Une transcription qui lisse un mot douteux a détruit la seule chose pour
% laquelle on la faisait — un lecteur doit pouvoir savoir, à chaque endroit,
% ce qui était sur le feuillet et ce qui a été ajouté. Les six macros qui
% suivent sont donc l'appareil critique, et non de la décoration.
%
% Les mêmes commandes sont reconnues par `scripts/render.mjs`, qui produit la
% vue de lecture du site. Une macro ajoutée ici doit l'être aussi là-bas,
% faute de quoi le rendu échouera bruyamment — ce qui est voulu : un rendu
% silencieusement faux est pire qu'un rendu absent.

\ProvidesPackage{grothendieck}[2026/08/08 Transcriptions du fonds Grothendieck]

\usepackage{amsmath,amssymb,amsthm}
\usepackage{mathrsfs}
\usepackage{stmaryrd}
% Les diagrammes commutatifs sont partout dans ce fonds : c'est la langue
% même de la géométrie algébrique de Grothendieck, et une transcription qui
% les rendrait en prose ne transcrirait rien.
\usepackage{tikz-cd}
% Français par défaut — c'est la langue des feuillets — mais l'anglais est
% chargé aussi : la traduction et le résumé sont des documents anglais, et la
% typographie française (espaces avant les deux-points) y serait une faute.
% Ces documents basculent par \selectlanguage{english} en tête de corps.
\usepackage[english,french]{babel}
\usepackage{fontspec}
\usepackage{geometry}
\geometry{a4paper,margin=25mm,marginparwidth=28mm}
\usepackage{marginnote}
\usepackage{xcolor}
% ulem, not soul. `soul`'s \st and \ul are built on a character-by-character
% scan that XeLaTeX's font machinery breaks: the document fails with an
% undefined control sequence rather than degrading. ulem does the same two jobs
% and is Unicode-engine safe. `normalem` stops it hijacking \emph, which the
% transcriptions use for Grothendieck's own emphasis.
\usepackage[normalem]{ulem}
\usepackage{hyperref}
\hypersetup{colorlinks=true,linkcolor=black,urlcolor=black,pdfborder={0 0 0}}

% --- Métadonnées du lot -----------------------------------------------------
% Reprises telles quelles par le script de rendu, qui les affiche en tête de
% la vue de lecture. Elles fixent ce que le fichier prétend couvrir : sans
% elles, une transcription flotte sans point d'attache dans un fonds de seize
% mille feuillets.

\newcommand{\folder}[1]{\gdef\grothFolder{#1}}
\newcommand{\batch}[1]{\gdef\grothBatch{#1}}
\newcommand{\leaves}[2]{\gdef\grothFirst{#1}\gdef\grothLast{#2}}
\newcommand{\foldertitle}[1]{\gdef\grothFoldertitle{#1}}
\gdef\grothFolder{?}\gdef\grothBatch{?}\gdef\grothFirst{?}\gdef\grothLast{?}\gdef\grothFoldertitle{}

% --- Le résumé --------------------------------------------------------------
%
% Il ouvre la lecture modernisée, sous le titre « Résumé », et il est composé
% en petit. Ce n'est pas un ornement : c'est le seul passage du document qui
% ne soit pas des mathématiques, et un lecteur doit pouvoir le reconnaître
% avant de l'avoir lu — puis le sauter s'il connaît déjà le sujet, ou s'y
% arrêter s'il ne le connaît pas. Le filet qui le suit marque la reprise du
% corps.

\newenvironment{resume}
  {\small\setlength{\parskip}{0.45em}}
  {\par\medskip\hrule\medskip}

% --- Mots-clés ---------------------------------------------------------------
%
% En anglais, en clôture du résumé de la lecture modernisée : le vocabulaire
% moderne sous lequel chercher ce que les feuillets construisent. Le manifeste
% du site (`npm run manifest`) les extrait pour étiqueter la cote entière —
% cette ligne est la source unique des tags, il n'y a pas de fichier à côté.
\newcommand{\keywords}[1]{\par\smallskip\noindent
  {\footnotesize\textsc{Keywords} --- \textit{#1}}\par}

% --- L'appareil critique ----------------------------------------------------

% Le numéro de feuillet, en marge. C'est l'ancre qui permet de régler un
% désaccord sur une formule en regardant *un* feuillet plutôt qu'une chemise
% de six cents. Le site s'en sert aussi pour tourner les pages du fac-similé
% quand on fait défiler la transcription.
\newcommand{\leaf}[1]{%
  \marginnote{\footnotesize\color{gray!70}\textsf{#1}}%
  \label{leaf:#1}\ignorespaces}

% Illisible : on ne devine pas. Un mot inventé qui se lit comme les autres est
% le pire résultat possible.
\newcommand{\ill}{\textcolor{red!60!black}{[\,\ldots\,]}}

% Lecture douteuse : le mot est donné, et signalé comme incertain.
\newcommand{\uncertain}[1]{\uline{#1}}

% Ajout de l'éditeur — une lettre manquante, un mot sous-entendu.
\newcommand{\add}[1]{\textcolor{blue!50!black}{[#1]}}

% Biffé par Grothendieck lui-même. Ce qu'il a rayé fait partie du document :
% une rature dit souvent où la pensée a changé de direction.
\newcommand{\struck}[1]{\sout{#1}}

% Note du transcripteur, sur l'état matériel du feuillet.
\newcommand{\note}[1]{\par\smallskip{\small\color{gray!60!black}\textit{[#1]}}\par\smallskip}

% Note marginale de Grothendieck, distincte du corps du texte.
\newcommand{\marginal}[1]{\par\smallskip\noindent\hspace{1em}%
  {\small\color{gray!60!black}#1}\par\smallskip}

% --- Titre ------------------------------------------------------------------

\AtBeginDocument{%
  \begin{center}
    {\large\bfseries Fonds Alexandre Grothendieck --- cote n\textsuperscript{o}~\grothFolder}\\[2pt]
    {\itshape\grothFoldertitle}\\[4pt]
    {\small feuillets \grothFirst--\grothLast\ (lot \grothBatch)}\\[2pt]
    {\footnotesize\color{gray!60!black}Transcription établie d'après le fac-similé numérisé
     par l'Université de Montpellier.\\
     Les lectures incertaines sont soulignées, les illisibles notées [\,\ldots\,],
     les ajouts entre crochets.}
  \end{center}
  \vspace{1em}}

\endinput


\folder{155}
\batch{1}
\pages{1}{10}
\dating{[à partir de 1984]}
\watermark{Édition de démonstration}
\foldertitle{[Rigidité de catégories tensorielles ?] : notes manuscrites (s.d.)}

\begin{document}

\note{les cinq feuillets de mathématiques du dossier portent sa pagination
propre, un chiffre cerclé en haut à droite, 1 à 5, sur les pages 1, 3, 5, 7 et
9 des archivistes. Le titre entre crochets du fonds est
celui des archivistes ; rien sur les feuillets ne porte de titre}

\page{1}

Soient $k$ un corps, $k'/k$ ext. finie.

La donnée d'une $\otimes$-catégorie \uncertain{ACU} $C$ sur $k$, et d'un
\struck{foncteur fibre (exact) \ill{} fidèle} $\otimes$-foncteur $k$-lin.
\uncertain{ACU} \add{exact et} \textbf{fidèle}
\[
F : C \longrightarrow \operatorname{Mod}_{f}(k')
\]
équivaut à la donnée

\note{l'abréviation qui qualifie ici la catégorie puis le foncteur est formée
d'un $A$ suivi de deux lettres liées ; elle est lue \emph{ACU} — associatif,
commutatif, unité — et n'est développée nulle part dans le dossier. La ligne
biffée est remplacée par l'insertion « exact et » au-dessus et par « fidèle »
au-dessous}

1°) d'un schéma affine $G'$ sur $k'$, muni d'une section
$\varepsilon$, i.e. d'une $k'$-alg. commutative $U$, \struck{\ill{}} \ill{}
$k'$-augmentation $\varepsilon : U \to k'$,

2°) sur le pro-objet \struck{$A$} $A = \check{U}$ \struck{(de
$\operatorname{Mod}_{f}(k')$)}, considéré comme pro-objet de
$\operatorname{Mod}_{f}(k')$, une structure de \struck{\ill{}}
\uncertain{commutative et associative}, admettant $\varepsilon$ comme unité,

\note{au-dessus de $A = \check{U}$, en interligne : $\overset{\text{déf}}{=}
\operatorname{Hom}_{k'}(U,k')$}

avec les conditions suivantes :

\begin{enumerate}
\item[(i)] l'homomorphisme $k'$-linéaire $\lambda \mapsto \lambda\varepsilon$,
$k' \to A$ est \textbf{multiplicatif} \note{en interligne, d'une main plus
rapide : « \ill{} la structure de pro-anneau de $A$ pour le \struck{\ill{}}
\ill{} de $k'$ »}

\item[(ii)] l'hom. \ill{}
\[
A \longrightarrow A \mathbin{\hat{\otimes}}_{k'} A
\]
déduit de la multiplication dans $U$ est « multiplicatif », i.e. prenant ses
valeurs dans
\[
\widetilde{\mathcal{J}}/\mathcal{J} \subset A \mathbin{\hat{\otimes}}_{k} A
\qquad
\bigl(\ \mathcal{J} = J\cdot(A \mathbin{\hat{\otimes}}_{k} A),\ \text{idéal} \ \bigr)
\]
\end{enumerate}

\note{sous $A \mathbin{\hat{\otimes}}_{k} A$, accoladé : « pro-$k$-alg.
prod\struck{uit} », et au-dessous « \struck{tens.} »}

\ill{} dont \uncertain{comp.} par
\[
J = \operatorname{Ker}(k' \otimes_{k} k' \to k'),
\qquad
\widetilde{\mathcal{J}} = \{\, \lambda \in A \mathbin{\hat{\otimes}}_{k} A
\mid \lambda\mathcal{J} \subset \mathcal{J} \,\}
\]

\note{les deux $J$ ne se distinguent sur le feuillet que par le tilde : $J$ est
le noyau de $k' \otimes_{k} k' \to k'$, $\mathcal{J}$ l'idéal qu'il engendre
dans $A \mathbin{\hat{\otimes}}_{k} A$, $\widetilde{\mathcal{J}}$ son
idéalisateur. La distinction typographique est de l'éditeur}

et à ce titre est \struck{\ill{}} \ill{} hom. d'anneaux $A \to
\widetilde{\mathcal{J}}/\mathcal{J}$.

\page{3}

\emph{Ex.} \quad $C = \operatorname{Mod}_{f}(k)$, \quad $F =$ le foncteur $M
\otimes_{k} k'$.

Alors $U = k' \otimes_{k} k'$ (structure de $k'$-alg.), section canonique
$\varepsilon : k' \otimes k' \to k'$, provenant du foncteur \ill{}

donc
\[
\check{U}^{(k')} = A \;\simeq\; k' \otimes_{k} \check{k}'^{(k)}
\;\simeq\; \operatorname{End}_{k}(k'),
\]
et \ill{} y prend les structures d'anneaux \ill{} habituelles.

\marginal{N.B. si $k'/k$ est \ill{} donc $U \simeq A^{\vee(k)}$ \ill{}
$U \xrightarrow{\ \Delta_{U}\ } U \otimes_{k} U$ transposé de $A
\otimes_{k} A \to A$ \ill{} \struck{\ill{}} \ill{} associative \ill{}
($\Delta_{U}(1) \neq 1$ en général), ($\otimes$ \ill{} injectif \ldots) et du
\ill{}}
\note{ce bloc court en biais, à environ $45^{\circ}$, le long du bord gauche du
feuillet, d'une main plus rapide que le corps ; il est donné aussi loin qu'il
se laisse lire}

$\langle$ On cherche conditions \add{intrinsèques} sur $C$ \add{le $\otimes$
est} qui assurent \ill{} pour \ill{} $U \subset$ \struck{\ill{}} $k'/k$ \ill{}
et $F$ comme dessus, \add{que} le foncteur des $A_{k}$-modules de t.f.
\note{l'indice de $A$ se lit $k$ ou $k'$ ; au-dessus de la lettre, en
interligne : « $= A \otimes_{k} k'$ »}
\[
M' \longrightarrow M' \otimes_{k' \otimes_{k} k'} k'
\]
soit \textbf{exact} (OK si $k'/k$ \ill{} \uncertain{étale}) et
\textbf{fidèle}.

\note{la ligne est ouverte par un crochet et séparée de ce qui précède par un
long trait courant d'un bord à l'autre ; les deux insertions « l'intrinsèque »
et « le $\otimes$ est » sont portées au-dessus de la ligne}

Suffit-il que $C$ soit \textbf{rigide}, et que $k \overset{\sim}{\to}
\operatorname{End}_{C}(1_{C})$ ? ? ?

\textbf{Exprimer la rigidité de $C$ !} \note{doublement souligné, dernière
ligne du feuillet}

\section*{Rigidité}

\note{titre de sa main, souligné, en tête du feuillet portant son chiffre 3}

\page{5}

$M \to \check{M}$ \struck{\ill{} $D(M)$} compatible avec foncteur fibre $F'$,
foncteur équivalence de catégories, donc \textbf{exact} donc cont. à droite,
\struck{\ill{}} donc donné par
\[
M \longmapsto \operatorname{Hom}_{A}(M, I)
\]
$I$ objet de $C$, muni d'opération de $A$ : $g$.
\[
I \;=\; \struck{\ill{}}\, D(A) \;\simeq\; \check{A}_{g}^{(k')} \;=\; U
\]

\note{sous le signe $\simeq$ : « compatibilité avec $F'$ »}

avec opération à g. de $A$ provenant par transposition \struck{des opérations}
de $A$ à droite sur $A_{g}$, dans $U$ \ill{} une $A$-module.

Il faut de faire opérer $A$ à gauche sur \ill{} cet objet de
$\operatorname{Mod}(A)$ \ill{} mettre sur $U$ une structure de
\struck{bimodule} module sur
\[
A \otimes A
\]
\note{sous la formule, entre parenthèses : « \ill{} opération déjà connue
\ill{} »}

Et les endomorphismes \ill{}

\marginal{\ill{} on \ill{} \ill{} $U$ \ill{} $D_{k'}(A)$ \ill{} $U$ \ill{}
$D_{k'}(A)$ \ill{} si $k'/k$ \ill{} fini \ill{}}
\note{second bloc en biais le long du bord gauche, même main rapide que celui
du feuillet portant son chiffre 2 ; seuls les deux $D_{k'}(A)$ et la mention de la
finitude de $k'/k$ s'en laissent tirer}

\ill{} (autrement, ils ne font pas voir \ill{} la $k'$-linéarité \ill{} pour
les $k'$-structures des deux côtés sur $U$ !)

$U$ \ill{} tel que \ill{} $A$-modules \ill{} un \ill{} \ill{} endomorphismes de
$A$ \ill{} $A$ module : autre \ill{}
\[
\operatorname{End}_{A}(U) \;\simeq\; \operatorname{End}_{A}(A_{d})^{\circ}
\;\simeq\; A^{\circ}
\]
donc il faut un \add{donné} \ill{}
\[
A \longrightarrow A^{\circ}
\]
\note{au-dessus de la flèche : « $a : a \mapsto \check{a}$ »}

\page{7}

N.B.
\[
D(M) \;=\; \operatorname{Hom}_{A}(M, U_{g})
\;=\; \operatorname{Hom}_{A}\bigl(M_{g}\,;\, \operatorname{Hom}_{k'}(A_{g}, k')\bigr)
\]
$\simeq$ \struck{$A_{d}$} modules des formes $k'$-bil.
\[
\varphi : A_{g} \times M_{g} \longrightarrow k'
\]
telles que
\begin{align*}
\varphi(\lambda u, x) &= \lambda\,\varphi(u,x), \\
\varphi(u \cdot v, x) &= \varphi(u, v.x),
\end{align*}
\note{en regard, la légende des variables : $\lambda \in k'$, $u \in
\uncertain{A}$, $x \in \uncertain{M}$}

i.e. compatible aux \ill{}, \uncertain{donne}
\[
A_{d} \otimes_{A} M \longrightarrow k'
\]
qui soit $k'$-linéaire, pour la structure de $k'$-module à gauche provenant de
$A_{d}$.

Or
\[
A_{d} \otimes_{A} M \;\simeq\; M \qquad \text{(comme du $k'$-module)}
\]
\[
1 \otimes x \longleftarrow x,
\qquad
\lambda(1 \otimes x) = \lambda \otimes x = 1 \otimes \lambda x \longleftarrow \lambda x
\]
donc en toutes fins
\[
F'(D(M)) \;=\; \check{M}^{(k')}
\]
et la donnée des $a : A \to A^{\circ}$ \ill{} \ill{} sur $\check{M}^{(k')}$ une
structure de $A$-module.

Il faut encore que
\[
\underbrace{F'(M) \otimes F'(D(M))}_{(M^{k'})^{\vee}}
\longrightarrow k' = F'(1_{C})
\]
\note{dans $F'(D(M))$ une première lettre est repassée et biffée ; au-dessus de
la ligne, biffée : $\check{M} \otimes \check{M}^{(k')}$ ; au-dessous, un second
essai biffé d'où l'on ne tire que $M^{k'}$}

provienne bien \struck{$A$} d'un hom.
\[
M \otimes D(M) \longrightarrow 1_{C}
\]

\page{9}

i.e. que
\[
M^{k'} \otimes_{k'} (M^{k'})^{\vee} \longrightarrow k'
\]
est $A$-linéaire, pour la structure \add{$A$} de gauche.

\note{le $A$ est inséré au-dessus du mot « structure »}

Il suffit \uncertain{en le prouver} pour
\[
M = A
\]
\[
A^{k'} \otimes_{k'} U \xrightarrow{\ \text{acc. can.}\ } k'
\]

$A \otimes$ \ill{}

\note{la page s'arrête là, sur une ligne commencée et abandonnée ; le reste du
feuillet est blanc, et c'est le dernier feuillet de mathématiques du dossier}

\end{document}
